\documentclass[a4paper,11pt]{article}

\usepackage[utf8]{inputenc}
\usepackage[T1]{fontenc}
\usepackage{geometry}
\usepackage{amsmath}
\usepackage{amssymb}
\usepackage{graphicx}
\usepackage{hyperref}
\usepackage{parskip}

\geometry{hmargin=2.5cm,vmargin=2.5cm}

\title{\textbf{Intelligent Neurons Network (INN)}\\ \large Towards Biologically Plausible Graph-Based Neural Processing}
\author{Antoine Salomon}
\date{\today}

\begin{document}

\maketitle

\begin{abstract}
Standard artificial neural networks (ANNs) rely on static, dense layers of simple perceptrons. While effective, this structure diverges significantly from biological brains, where neurons are specialized, state-dependent, and interact through dynamic pathways. This paper proposes a novel architecture: the \textbf{Intelligent Neuron (IN)}. By endowing each neuron with memory (LSTM), a communication interface (Key/Query attention), and a specific role within a graph topology, we introduce the \textbf{Intelligent Neurons Network (INN)}. This architecture aims to bridge the gap between deep learning and biological functional specialization.
\end{abstract}

\section{Introduction}

The human brain is not a stack of matrix multiplications. It is a complex graph of specialized cells. Sensory neurons (retina, cochlea) encode raw signals, transmitting them to specialized cortical columns that process information before sending commands to motor neurons (muscles).

Current Deep Learning models, such as Transformers, rely on a monolithic approach: broad layers processing information uniformly. We propose a shift towards a more granular and modular approach. The goal is to mimic the biological functional chain:
\begin{itemize}
    \item \textbf{Input Neurons:} Analogous to sensory organs (eyes, ears), receiving external data.
    \item \textbf{Transmission Neurons:} Analogous to the cortex, processing and routing information.
    \item \textbf{Action Neurons:} Analogous to motor outputs, producing the final system response.
\end{itemize}

In this paradigm, the network is no longer a fixed circuit, but a dynamic society of agents (neurons) communicating via a content-addressable bus system.

\section{Why INN is Novel and Promising}

The INN architecture introduces three key innovations compared to traditional RNNs or Transformers:

\begin{enumerate}
    \item \textbf{Graph Topology Learning:} Unlike dense layers where connectivity is fixed, INN learns \textit{who speaks to whom}. The topology emerges during training through attention mechanisms.
    \item \textbf{Functional Specialization:} Because communication is selective (via Keys and Queries), neurons are encouraged to specialize. One neuron might become a "punctuation detector" listened to by "capitalization neurons," creating explicit functional sub-circuits.
    \item \textbf{Asynchronous Complexity:} The graph structure allows for varied processing depths. Information can flow directly from Input to Action for reflexes, or loop through multiple Transmission neurons for complex reasoning, mimicking "System 1 vs System 2" thinking.
\end{enumerate}

\section{Intelligent Neurons (The Micro-Architecture)}

An Intelligent Neuron is not a scalar scalar activation unit. It is a vector processing unit with three components:

\subsection{1. Signal Receiver (The Ear)}
The neuron filters incoming information using a \textbf{Query Vector} ($Q$). It does not passively receive all inputs; it actively "listens" for specific signal patterns relevant to its function.
\[ \text{Input}_{\text{filtered}} = \text{ReceiverFF}(\sum_{j} \alpha_{ij} \cdot \text{Signal}_j) \]
where $\alpha_{ij}$ is the attention weight determined by the match between the neuron's Query and the sender's Key.

\subsection{2. Core Memory (The Brain)}
The processing heart is a Long Short-Term Memory (LSTM) cell. It maintains a hidden state ($h_t$) and a cell state ($c_t$), allowing the neuron to have temporal context.
\[ h_t, c_t = \text{LSTM}(\text{Input}_{\text{filtered}}, h_{t-1}, c_{t-1}) \]
This memory allows the neuron to be state-dependent: its reaction depends not just on current input, but on its history.

\subsection{3. Signal Sender (The Voice)}
The neuron broadcasts its processed state to the network. It projects its state into a \textbf{Signal Vector} and tags it with a \textbf{Key Vector} ($K$).
\[ \text{Output} = \text{SenderFF}(h_t) \]
The Key ($K$) acts as a metadata tag (e.g., "I am broadcasting a verb"), allowing other neurons to decide whether to listen.

\section{Intelligent Neurons Network (The Macro-Architecture)}

The INN is a connected graph of these Intelligent Neurons. The "glue" holding them together is the Global Attention Mechanism.

At each time step $t$:
\begin{enumerate}
    \item \textbf{Broadcast:} Every neuron exposes its Key ($K$) and Query ($Q$).
    \item \textbf{Matchmaking:} A global attention matrix $A$ is computed:
    \[ A_{ij} = \text{softmax}\left(\frac{Q_i \cdot K_j^T}{\sqrt{d_k}}\right) \]
    This determines the bandwidth of the connection from Neuron $j$ to Neuron $i$.
    \item \textbf{Transmission:} Signals are routed according to $A$.
    \item \textbf{Update:} All neurons update their internal LSTM state simultaneously.
\end{enumerate}

This structure is fully differentiable, allowing end-to-end training via Backpropagation Through Time (BPTT).

\section{Vectorized \& Hybrid Implementation}

To scale the architecture to modern GPU capabilities (A100), we transitioned from an object-oriented approach to a fully vectorized implementation.
The \textbf{Hybrid INN} introduces two populations of neurons within the same network:
\begin{itemize}
    \item \textbf{Static Neurons ($N_{static}$):} These neurons have learnable but fixed $K$ and $Q$ parameters. They form the "backbone" of the network, capturing stable rules (grammar, syntax).
    \item \textbf{Dynamic Neurons ($N_{dynamic}$):} These neurons generate $K_t$ and $Q_t$ at each time step via a projection of their hidden state $h_t$. They act as "liquid resources," recruited on-demand for complex context switching or algorithmic tasks.
\end{itemize}

The global attention matrix $A$ is constructed by concatenating static and dynamic keys/queries:
\[ Q_{total} = [Q_{static} \parallel Q_{dynamic}(h_t)] \]
This allows the network to seamlessly mix reflex pathways and adaptive reasoning.

\section{Experimental Validation: Multi-Tasking Plasticity}

To test the hypothesis of functional specialization, we designed a "Dual-Mode" experiment combining two radically different cognitive tasks:
\begin{enumerate}
    \item \textbf{Literary Generation:} Shakespearean text (Context-heavy, ambiguous syntax).
    \item \textbf{Arithmetic Logic:} Symbolic additions in the format \texttt{MATH: A+B=C} (Rule-based, strict syntax).
\end{enumerate}

The dataset alternates between these two modes randomly, forcing the network to "context-switch" constantly. We trained a Hybrid INN (40 Static, 20 Dynamic Neurons) for 30 epochs.

\subsection{Results}

\subsubsection{Convergence & Performance}
Contrary to standard Multi-Task learning where interference often degrades performance, the INN trained on the mixed dataset achieved a \textbf{lower Loss (1.84)} than the INN trained on Shakespeare alone (1.99).
This suggests that the structural constraints of arithmetic helped the network organize its latent space more efficiently, benefitting the literary task as a side effect.

\subsubsection{Emergent Arithmetic}
The network successfully learned the syntax of equations and the approximate magnitude of results after only 20 minutes of training:
\begin{verbatim}
Generated: MATH: 378+301=681  (Ground Truth: 679)
Generated: MATH: 729+607=1241 (Ground Truth: 1336)
\end{verbatim}
While not perfectly precise (which would require Curriculum Learning), the model captured the algorithmic nature of the task (column-wise addition logic) without any explicit programming.

\subsection{Visualizing Neuroplasticity}

The evolution of the Attention Matrix reveals the network's self-organization strategy:

\begin{figure}[h]
    \centering
    % Placeholder for your image: attention_epoch_10.png
    \includegraphics[width=0.8\textwidth]{attention_matrices/attention_epoch_10.png}
    \caption{\textbf{Exploration Phase (Epoch 10):} The network activates multiple Dynamic Neurons (TD columns) in a broad search for solution pathways. 5 distinct dynamic hubs are active.}
\end{figure}

\begin{figure}[h]
    \centering
    % Placeholder for your image: attention_epoch_20.png
    \includegraphics[width=0.8\textwidth]{attention_matrices/attention_epoch_20.png}
    \caption{\textbf{Consolidation Phase (Epoch 20):} Activity concentrates on 3 highly specialized Dynamic Neurons. The network has "pruned" unnecessary connections, automating the task processing.}
\end{figure}

This confirms our core hypothesis: \textbf{The INN does not just learn weights; it learns a topology.} It recruits dynamic resources when the task difficulty increases (switching to MATH mode) and releases them when stable patterns are found.

\section{Conclusion}

The Intelligent Neurons Network represents a step away from the "black box" paradigm of dense matrix multiplication. By reintroducing biological constraints---functional roles, sparse dynamic communication, and local memory---we created a system capable of:
\begin{itemize}
    \item \textbf{Self-wiring:} Learning its own graph structure.
    \item \textbf{Liquid Adaptation:} Changing its connectivity on the fly.
    \item \textbf{Efficient Multi-Tasking:} separating cognitive domains in orthogonal subspaces.
\end{itemize}

While currently slower to train than highly optimized Transformers, the INN offers a path towards \textbf{System 2 AI}: systems that can reason, adapt, and explain their internal routing, rather than just predicting the next token statistically.

\end{document}
